% *********************************************************************
% © 2016–2026 Jeremy Sylvestre
%
% Permission is granted to copy, distribute and/or modify this document
% under the terms of the GNU Free Documentation License, Version 1.3 or
% any later version published by the Free Software Foundation; with no
% Invariant Sections, no Front-Cover Texts, and no Back-Cover Texts. A
% copy of the license is included in the appendix entitled “GNU Free
% Documentation License” that appears in the output document of this
% PreTeXt source code. All trademarks™ are the registered® marks of
% their respective owners.
%
% *********************************************************************
\tdplotsetmaincoords{50}{25}
\begin{tikzpicture}[
	point/.style={circle,draw,very thin,fill,inner sep=0pt,minimum size=4pt},
	vector/.style={-latex},
	tdplot_main_coords
]
	% artificially widen image because it's in a sidebyside with a <p>
	% and want to give caption more room to expand
	\draw[black!5!white,very thin] (-1.35,-1.3,-0.05) to (-1.35,-1.3,-0.25) to (-1.25,-1.25,-0.25);
	\draw[black!5!white,very thin] (2.1,1.9,1.2) to (2.1,1.9,1.35) to (2,1.85,1.35);


	%draw the axes
	\draw[->] (-0.5,0,0) to (2,0,0) node[right]{$x$};
	\draw[->] (0,-0.5,0) to (0,2,0) node[right]{$y$};
	\draw[->] (0,0,-0.5) to (0,0,2) node[right]{$z$};
	\node[point] at (0,0,0) (o) [label=below left:{$(0,0,0)$}] {};
	\node[point] at (1,0,0) (e1) [label=above right:{$(1,0,0)$}] {};
	\node[point] at (0,1,0) (e2) [label=right:{$(0,1,0)$}] {};
	\node[point] at (0,0,1) (e3) [label=left:{$(0,0,1)$}] {};
	\draw[vector,very thick] (o) to node[below] {$\uvec{e}_1$} (e1);
	\draw[vector,very thick] (xyz cs:) to node[right] {$\uvec{e}_2$} (e2);
	\draw[vector,very thick] (xyz cs:) to node[left] {$\uvec{e}_3$} (e3);
\end{tikzpicture}
